\section{Methodology}

\subsection{Research Methodology}

% Describe the overall research approach and methodology used in the study.


\subsection{Experimental Methodology}
This chapter describes how the experimental work in this thesis was designed, conducted, and evaluated. It covers the design of experiments, the test setup and equipment used and the procedure followed during testing.

% Introduce the experimental methodology and objectives.

\subsubsection{Design of Experiments}
The experiments in this thesis are designed to evaluate how different pipeline configurations affect end-to-end streaming latency, video quality and resource usage. To do this in a structured way, the work draws on the principles of Design of Experiments (DoE), which is a formal framework for planning tests so that the results can be interpreted clearly and with confidence \cite{montgomery2017}.

The experimental design follows the One Factor At a Time (OFAT) strategy. This means one variable is changed between configurations while all others are kept at a fixed baseline value. Montgomery describes OFAT as a well established approach when the goal is to understand the individual contribution of each factor separately, rather than exploring how factors interact with each other \cite{montgomery2017}.

The factors tested in this thesis are codec, resolution, transport protocol, encryption, bandwidth condition and encoding method. While GStreamer allows a large number of additional parameters to be adjusted, such as buffer sizes, queue configurations, bitrate control modes and keyframe intervals, expanding the test matrix further was not feasible within the time constraints of a single master thesis. The chosen parameters are considered sufficient to identify the most significant performance differences relevant to the use case.

The outcomes measured for each configuration are end-to-end latency, CPU usage and visual quality. Packet loss and throughput are also recorded for configurations that involve simulated poor network conditions. All configurations are described in the following section.
% Describe the experimental design, variables, and planning.

\paragraph{Table}
% Insert experiment design table here.

\subsubsection{Test Setup and Equipment}

% Describe the hardware, software, instruments, and setup used.

\paragraph{Table}
% Insert equipment/setup table here.

\subsubsection{Test Procedure}

% Explain the step-by-step testing procedure.

\subsubsection{Controlling for Unmeasured and Hidden Variables}

% Explain how uncontrolled or hidden variables were minimized.

\subsection{Measurement Methodology}

% Describe the measurement methodology used in the project.

\subsubsection{How You Used That Design During Testing}

% Explain the practical use of the experimental design during testing.
% Include:
% - Testing procedure
% - Accuracy considerations
% - Error handling

\subsubsection{What Was Measured}

% Describe all measured variables and parameters.

\subsubsection{Measurement Method}

% Explain how measurements were collected and recorded.

\subsubsection{Accuracy and Error Sources}

% Discuss measurement accuracy and possible error sources.

\subsubsection{Data Recording and Processing}

% Explain how data was stored, processed, and analyzed.

\subsection{Statistical Analysis}

% Describe statistical methods and analysis techniques.

\subsection{Limitations}

% Discuss limitations of the methodology and experiments.

\subsection{Reliability}

% Explain how reliability and repeatability were ensured.

\subsection{Validity}

% Explain validity of the measurements and conclusions.


\subsection{Validity}

\subsection{Reliability}

\begin{comment}
\subsection{Research Design}

This thesis follows an applied research approach, combining Design Science Research (DSR) and experimental research to address both the development and evaluation aspects of the work. This combination is well suited for engineering projects where the goal is to produce a functional artifact while also generating knowledge through systematic testing and measurement.

The DSR component governs the design and development of the video streaming pipeline. Following the principles outlined by Hevner et al. [source], the research contribution is embedded in the artifact itself, meaning the pipeline design, the architectural decisions made, and the reasoning behind them all form part of the research output. The problem is clearly grounded in practice, the solution is built and refined iteratively, and the artifact is evaluated against the requirements defined in the previous chapter.

The experimental research component covers the testing and evaluation phase. Drawing on the framework described by Wohlin et al. [source], the pipeline is tested across a range of configurations in a controlled setting. Variables such as codec choice, payload settings, and network conditions are systematically varied, and end-to-end latency is measured and compared across configurations. This allows conclusions to be drawn about which configurations perform best and what trade-offs exist between latency and video quality.

Together these two approaches reflect the dual nature of the thesis: it is both a development project and a measurement study, and the methodology is designed to support both goals.

\end{comment}

\subsection{Experimental Methodology}
\subsubsection{Measurement Method}

